% ============================================================================
%  Sorting Algorithms: Overview and Objectives  (CIS 4020)
%  Requires classnotes.sty in the same folder or your local texmf tree.
%  Compile twice with pdflatex so "Page X of Y" resolves.
% ============================================================================
\documentclass[11pt]{article}
\usepackage[margin=1in]{geometry}
\usepackage[pagenumbers]{classnotes}
\usepackage{needspace}

\classnotesmeta{Sorting Algorithms: Overview and Objectives (CIS 4020)}

\begin{document}

\classnotestitle{CIS 4020 \quad$\cdot$\quad DATA STRUCTURES}{Sorting Algorithms: Overview and Objectives}{Bubble, selection, insertion, merge, and quick sort}

\section{Overview}

Sorting means arranging a collection of items into a specified order, such as smallest to
largest. It is one of the oldest and most studied problems in computer science, and it
appears constantly in real programs: ordered search results, ranked lists, leaderboards, and
data prepared so that a fast search can work at all. The binary search from the algorithm
analysis unit only works on sorted data, so sorting is what makes that speedup available.

This unit covers five sorting algorithms. Three are simple and closely related:
\textbf{bubble sort}, \textbf{selection sort}, and \textbf{insertion sort}. Two are
divide and conquer algorithms: \textbf{merge sort} and \textbf{quicksort}. Libraries already
provide a sort, so it is fair to ask why we write our own. The answer is that each of these
five makes a different trade-off between speed, memory, and behavior on different kinds of
input, and working through them is the best practice available for reasoning about an
algorithm rather than only using one.

This unit is also where two earlier units get used together. Recursion gives merge sort and
quicksort their structure, and algorithm analysis gives us the tools to count the work each
sort does and compare them fairly.

For every sort, the goal is the same four things: you can \textbf{perform} it by hand,
\textbf{code} it, \textbf{analyze} it, and \textbf{compare} it with the others. The larger
goal is the method itself. Once you can carry those steps through five sorts, you can carry
them through a sort you have never seen. By the end of algorithm analysis and sorting
together, you should also be able to turn the same method on your own solution to a problem:
count what it does, name its growth rate, and decide whether it is good enough.

As in earlier units, we work in \textbf{two languages at once}. Computer Science students
write \textbf{C/C++}; CIS students write \textbf{Java}. Each of the five sorts is the same
algorithm in both languages, and only the syntax and the way an array is handed to a method
differ. Where those differences matter, we call them out.

The most useful preparation is to trace each sort by hand on a small array of five or six
numbers before you look at any code. If you can predict the array after every pass, the code
is much easier to write and to debug.

\Needspace{12\baselineskip}
\section{Learning objectives}

By the end of this unit you should be able to:

\subsection*{Remember and understand}
\begin{itemize}[leftmargin=1.4em]
  \item Define sorting, and state what it means for an array to be in ascending or
    descending order, including how equal values are treated.
  \item Describe in plain language the central idea of each of the five sorts: bubble sort
    swaps neighboring items that are out of order, selection sort repeatedly finds the
    largest item in the unsorted part of the array and swaps it into the last unsorted
    position, insertion sort grows a sorted region one item at a time, merge
    sort splits the array and then merges the sorted halves, and quicksort partitions the
    array around a pivot.
  \item Define the vocabulary of sorting: pass, comparison, swap, stable, in place, pivot,
    partition, and merge.
  \item State the best-case, average-case, and worst-case Big-O class for each of the five
    sorts.
  \item Explain, conceptually, what the pivot does in quicksort and why the choice of pivot
    affects how well the algorithm performs.
  \item Recognize merge sort and quicksort as divide and conquer algorithms, and name the
    divide, conquer, and combine parts of each.
\end{itemize}

\subsection*{Apply}
\begin{itemize}[leftmargin=1.4em]
  \item Trace each of the five sorts by hand on a small array, showing the contents of the
    array after each pass, split, merge, or partition.
  \item Write a working implementation of each sort in Java (CIS) or C/C++ (CS), using loops
    for the three simple sorts and recursion for merge sort and quicksort.
  \item Count the comparisons and swaps a given sort performs on a specific small input, and
    check that count against your trace.
  \item Apply the counting method from algorithm analysis to a sort's loops to find its
    Big-O class. For example, a loop nested inside a loop over the same array points to
    O($n^2$).
\end{itemize}

\subsection*{Analyze}
\begin{itemize}[leftmargin=1.4em]
  \item Determine which input produces the best case and which produces the worst case for
    each sort (for example, an already sorted array, a reverse-ordered array, or a random
    one), and explain why these are not the same for every sort.
  \item Explain why bubble sort, selection sort, and insertion sort are O($n^2$) by
    examining how their loops are structured, and identify which of the three does less work
    on an already sorted array, and why.
  \item Explain why merge sort is O($n \log n$): the array is halved about $\log n$ times,
    and the merging at each level does about $n$ work in total. Connect the halving to the
    binary search analysis.
  \item Explain why quicksort averages O($n \log n$) but degrades to O($n^2$) in its worst
    case, in terms of how evenly each partition splits the array.
  \item Explain why counting comparisons alone can be misleading, and distinguish a sort
    that does many comparisons but few swaps from one that does the opposite.
  \item Compare the five sorts on stability, whether they sort in place, and how much extra
    memory they need.
  \item Identify the base case, the recursive step, and the point where the real work
    happens in merge sort and in quicksort, and relate each to the recursion unit.
  \item Given a sort you have not seen before, such as a variation of one of the five or a
    new sort described in pseudocode, trace it, decide whether it is correct, and determine
    its Big-O class using the same steps you used for the five in this unit.
\end{itemize}

\subsection*{Evaluate}
\begin{itemize}[leftmargin=1.4em]
  \item Choose a sorting algorithm for a described situation (for example, a small array, a
    nearly sorted array, limited memory, a need for stability, or a need for a guaranteed
    worst case), and justify the choice against the alternatives.
  \item Critique a claim such as ``quicksort is always the fastest sort'' or ``this sort is
    O($n$) because it has only one loop,'' using best, average, and worst cases and the work
    hidden inside each loop.
  \item Judge whether a given sorting implementation is correct and efficient, and identify
    and explain any error, such as an off-by-one boundary, a missing base case, or
    unnecessary passes.
  \item Analyze and judge your own solution to a problem, sorting or otherwise: count what it
    does, identify its dominant term and its Big-O class, decide whether that is acceptable
    for the size of input you expect, and propose an improvement if it is not.
\end{itemize}

\end{document}
